%! Author = sriadityavedantam
%! Date = 3/25/26

\section{Discussion}

\lesson{41}{Monday 17 November 2025 9:10}{Day 41}

\begin{example}[The Fundamental Theorem of Calculus]
    \begin{gather*}
        \textbf{(i) } \int_a^b F\uptick(x)\,dx = F(b) - F(a).\\
        \textbf{(ii) } \text{If } G(x) = \int_a^x f(t)\,dt,\text{ then } G\uptick(x) = f(x).\\
    \end{gather*}
    This is a statement about the inverse relationship between differentiation and integration.
\end{example}

\begin{remark}
    The best way to think of integration is not only as the inverse process of differentiation.
    Historically, it was viewed this way by mathematicians such as Leibniz and Newton.
    However, the modern definition of integration is built from limits of sums.
    For example,
    \[
        h(x) = \begin{cases}
                   1,&\text{for } 0\leq x<1\\
                   2,&\text{for } 1\leq x\leq 2
        \end{cases}
    \]
    is Riemann integrable on $[0,2]$, even though it is not continuous at $x=1$.
    In calculus, we study Riemann sums, often drawn graphically as rectangular areas:
    \[
        \sum_{k=1}^n \underbrace{f(c_k)}_{\text{height}} \underbrace{(x_k - x_{k-1})}_{\text{width}}.
    \]
    As the partition gets finer, these sums approach the integral.
    Hence
    \[
        \lim_{n\to\infty} \sum_{k=1}^n f(c_k)(x_k - x_{k-1}) = \int_a^b f(x)\,dx,
    \]
    when the function is Riemann integrable.
\end{remark}